\section{OptiHarness}
\label{sec:method}

\subsection{Overview}

OptiHarness converts completed code-agent executions into updates of the harness used by later executions. Let the harness at cycle $c$ be
\begin{equation}
    h_c=(T_c,E_c,I_c),
\end{equation}
where $T_c$ is the model-visible tool registry, $E_c$ implements the registered tools, and $I_c$ gives generic cost-aware behavioral guidance. The base model and its parameters remain fixed. Given $h_c$ and a set of tasks sampled from $\taskdist$, a completed cycle yields a folded staged harness $g_c^\star$. If every batch is rejected or makes no change, then $g_c^\star=h_c$; an incomplete cycle produces no eligible handoff. The cycle also produces an audit trail containing annotated trajectories, focused proposal records, candidate diffs, and screening decisions.

Figure~\ref{fig:overview} shows the data flow. The current harness is fixed during rollout, so adaptation never rewrites messages already accumulated by the agent. Completed trajectories are annotated and reduced to bounded update hypotheses. A proposer edits the three coupled harness artifacts in staged state, and a static gate accepts or restores each batch. Only the final folded state is eligible for a later deployment.

\begin{figure*}[t]
\centering
\begin{tikzpicture}[
    >=Latex,
    box/.style={draw=black!50,rounded corners=2pt,align=center,
        font=\footnotesize,text width=37mm,minimum height=14mm,inner sep=2.5pt},
    rollout/.style={box,fill=gray!7},
    analysis/.style={box,draw=blue!60!black,fill=blue!7},
    proposal/.style={box,draw=orange!75!black,fill=orange!9},
    screen/.style={box,draw=green!45!black,fill=green!7},
    callout/.style={draw=red!55!black,rounded corners=2pt,fill=red!5,
        align=center,font=\scriptsize,text width=82mm,inner sep=3pt},
    flow/.style={->,thick,draw=black!70},
    loop/.style={->,thick,dashed,draw=orange!75!black}
]
\node[rollout] (s1) at (0,0) {\textbf{1. Current harness $h_c$}\\
$T$: tool registry\\$E$: dispatcher\\$I$: instruction};
\node[rollout] (s2) at (4.5,0) {\textbf{2. Fixed-harness rollout}\\
append-only actions and observations\\
unchanged-prefix reuse remains possible};
\node[analysis] (s3) at (9.0,0) {\textbf{3. Prefix-only annotation}\\
predicted dependencies, type, and state};
\node[analysis] (s4) at (9.0,-2.3) {\textbf{4. Focused evolving records}\\
observed $s^{-}\rightarrow$ proposed $s^{+}$\\
skip, merge, pivot, or bound output};
\node[proposal] (s5) at (4.5,-2.3) {\textbf{5. Unified staged update}\\
$\widetilde g=(\widetilde T,\widetilde E,\widetilde I)$};
\node[screen] (s6) at (0,-2.3) {\textbf{6. Static gate}\\
accept $g_b$ or restore $g_{b-1}$};

\draw[flow] (s1) -- (s2);
\draw[flow] (s2) -- (s3);
\draw[flow] (s3) -- (s4);
\draw[flow] (s4) -- (s5);
\draw[flow] (s5) -- (s6);
\draw[loop] (s6.south) to[out=-90,in=-90,looseness=1.15]
    node[below,font=\scriptsize] {next evidence batch} (s5.south);
\draw[flow] (s6.west) to[out=180,in=180,looseness=1.25]
    node[left,font=\scriptsize,align=center] {completed fold\\to next cycle} (s1.west);
\node[callout] at (4.5,-4.2)
    {The gate checks contracts and bounded semantic evidence; it does not execute the candidate or prove cost and utility improvement.};
\end{tikzpicture}
\caption{\textbf{Overview of \sys{}.} The harness comprises tool registry $T$,
dispatcher $E$, and behavioral instruction $I$ and remains fixed while a
rollout appends actions and observations. This leaves the accumulated history
unchanged and permits exact-prefix reuse when the provider recognizes that
prefix. A future-suffix-blind annotator
predicts action dependencies, types, and states. The resulting graph and
bounded behavioral patterns produce local records from an observed slice
$s^{-}$ to an unexecuted proposed slice $s^{+}$. These records guide a coupled
harness update. Deterministic contract checks and bounded semantic review
either accept the staged candidate or restore the preceding harness. Only a
completed fold is eligible for the next rollout.}
\label{fig:overview}
\end{figure*}


\subsection{Collecting trajectories under a fixed harness}

For a task $x\sim\taskdist$, the current harness produces the rollout object $\zeta_{h_c}(x)$ from Equation~\ref{eq:zeta} and its action projection
\begin{equation}
    \tau_{h_c}(x)=((a_1,o_1),\ldots,(a_n,o_n)).
\end{equation}
Here $a_i$ contains the tool calls made at action step $i$, and $o_i$ contains the corresponding environment results. The same $h_c$ is used throughout the rollout. New responses and observations are appended to the message sequence; the adaptation process starts only after the rollout finishes. Cycle~1 may use a pre-materialized baseline collected under the declared initial harness, while later cycles use the most recently handed-off harness.

Before a trajectory is admitted as adaptation evidence, OptiHarness checks execution completeness rather than benchmark success. Each requested task must correspond to one trial with a parseable result, a structurally valid trajectory, and at least one agent tool action. A zero verifier reward remains admissible because an unsuccessful but completed execution can reveal costly failure patterns. Missing, duplicate, exception-terminated, and action-free trials are excluded from the cycle.

\subsection{Future-suffix-blind action annotation}

The same shell command can be necessary in one context and redundant in another. OptiHarness therefore predicts dependencies between actions instead of assigning cost labels from surface form alone. For action $i$, the annotator receives the task, the raw action--observation steps $1$ through $i-1$, and the current pair $(a_i,o_i)$. It cannot observe any step after $i$. From this input it predicts
\begin{equation}
    \widehat P_i\subseteq\{0,1,\ldots,i-1\},\qquad
    u_i\in\{\texttt{read},\texttt{write},\texttt{verify},\texttt{explore}\},
\end{equation}
and
\begin{equation}
    v_i\in\{\texttt{success},\texttt{fail}\}.
\end{equation}
Index $0$ denotes the initial task state. The predecessor set $\widehat P_i$ records earlier actions predicted to be needed to generate action $i$; $u_i$ describes the operation type; and $v_i$ records whether the action produced a usable result. A rule based on return codes and explicit error markers supplies $v_i$ when the model output is malformed. The dependency prediction, operation type, and state otherwise come from the same annotation call.

The predictions define a directed acyclic graph
\begin{equation}
    \widehat G_\tau=(V_\tau,\widehat E_\tau),\qquad
    \widehat E_\tau=\{(j,i):j\in\widehat P_i,\ j>0\}.
\end{equation}
The graph is acyclic by construction because every predecessor index is smaller than its dependent action. Future-suffix blindness prevents a later successful action from being shown to the annotator when it explains an earlier decision. It does not make the predicted edges causal or correct. Step-level checkpointing changes only execution scheduling and recovery; it adds no scientific validation to the graph.

\subsection{Constructing focused proposal records}

\paragraph{Predicted final closure.}
Starting from the final action $n$, OptiHarness follows predicted predecessor edges backward and obtains
\begin{equation}
    \widehat\closure_\tau=\{n\}\cup\operatorname{Anc}_{\widehat G_\tau}(n).
\end{equation}
This set is minimal only under the predicted graph and the unchecked choice of the final action as anchor. We use it to localize possible waste, not as proof that the omitted actions are unnecessary to solve the task.

\paragraph{Graph-derived patterns.}
The first pattern is a contiguous region outside $\widehat\closure_\tau$. Such a region suggests that later instructions could discourage a dead-end exploration. The second pattern is a set of successful structured actions with the same predicted direct predecessors, operation type, and state. These sibling actions suggest a batched tool that can issue the operations together. Candidate write operations are merged only when their known target files are disjoint. In both cases, the builder retains a small amount of predecessor and consumer context around the target actions.

\paragraph{Behavioral and cost patterns.}
The graph alone does not expose every source of inference cost. OptiHarness also detects consecutive failed actions followed by a successful pivot, a dependency-relevant action with a long observation, and an operator family that recurs often enough to accumulate substantial output. These patterns respectively suggest an earlier strategy change, a bounded-output interface, or a reusable batched capability. When no primary pattern is found, a bounded phase with a common operation type supplies fallback evidence. These proposed transformations are design hypotheses; they are not replayed trajectories.

Each resulting record follows the object defined in Section~\ref{sec:problem}:
\begin{equation}
    e=(s^{-},s^{+},R_e,z,\pi_e).
\end{equation}
The observed slice $s^{-}$ contains the selected actions and their actual observations. The proposed slice $s^{+}$ removes, merges, or contracts the target pattern. The route $R_e\subseteq\{T,E,I\}$ states which harness artifacts the record should inform, $z$ identifies the detected pattern, and $\pi_e$ records its source and evidence status. To avoid repeating unchanged observations in the update prompt, OptiHarness serializes $s^{+}$ as a structural delta over $s^{-}$. It keeps a bounded number of high-scoring, type-diverse records from each source trajectory and removes any record containing actions that seek private answers, ground truth, or task-specific verifier data.

The implementation currently assigns graph-derived records the legacy status name \texttt{dependency\_validated}. This label overstates the evidence: neither the predicted graph nor the proposed slice has been validated by replay. Such records are treated in this paper as predicted-dependency hypotheses, and the legacy label must be renamed or rejected before a scientific evaluation uses the recorded provenance.

\subsection{Proposing a unified harness update}

All records in a cycle were collected under the same cycle-start harness $h_c$. They are grouped by intended route, detector type, and evidence status, then ordered round-robin across source tasks. Grouping gives a batch a coherent update target, while source rotation reduces the chance that a candidate is driven entirely by one task. If the prompt budget cannot hold a full batch, the batch is split. A record that cannot fit by itself stops processing instead of being silently dropped.

Let $g_{c,0}=h_c$ be the initial staged state and let $B_{c,b}$ denote update batch $b$. The proposer receives $g_{c,b-1}$, the records in $B_{c,b}$, and their observed-to-proposed deltas, and writes a candidate
\begin{equation}
    \widetilde g_{c,b}=\operatorname{Propose}(g_{c,b-1},B_{c,b}).
\end{equation}
The candidate may add, repair, merge, or remove generic tools, change their dispatch logic, and revise the behavioral instruction. It may also leave the harness unchanged when the evidence is weak or an existing capability already covers the pattern. The registry and dispatcher are presented together because a schema without matching execution semantics is unusable. Instructions are included in the same object because batching and stopping behavior cannot always be expressed through a tool interface.

The machine-enforced edit boundary is the three-artifact set $\{T,E,I\}$. A narrower $R_e$ is communicated to the proposer but is not enforced per record. This distinction matters when attributing a candidate change to its supporting evidence.

\subsection{Static screening and serial folding}

Each batch begins from an immutable copy of $g_{c,b-1}$ and is edited in a separate candidate state. A deterministic screen checks that the registry is valid JSON, tool names and parameter schemas are well formed, the dispatcher is valid Python with the required interface, registered tools have corresponding dispatch behavior, the instruction remains bounded, the diff remains reviewable, and no file outside $\{T,E,I\}$ changed. Failure of any check rejects the candidate without calling the semantic judge.

Candidates that pass these contracts receive a bounded-view semantic review. The judge sees the batch evidence, a summary of the complete candidate harness, and the candidate diff. It independently scores cross-task generality, prospective net API cost, correctness, internal consistency, and evidence alignment and minimality. Let $D(\widetilde g_{c,b})$ denote deterministic validity and let $S_k(\widetilde g_{c,b},B_{c,b})\in\{0,1,2,3,4\}$ be semantic score $k$. With threshold $\gamma=3$, the internal batch decision is
\begin{equation}
g_{c,b}=
\begin{cases}
\widetilde g_{c,b}, &
\substack{D(\widetilde g_{c,b})=1,\\
\min_k S_k(\widetilde g_{c,b},B_{c,b})\geq\gamma,\\
\text{and no blocking issue is reported},}\\
g_{c,b-1}, & \text{otherwise}.
\end{cases}
\label{eq:batch-fold}
\end{equation}
Transport errors, malformed judge responses, and below-threshold scores fail closed. Acceptance promotes only the three allowed artifacts from the candidate snapshot. Rejection restores the exact pre-batch state. Later batches see the latest accepted staged state, although their evidence still comes from rollouts under $h_c$.

This screening step evaluates contracts and prospective evidence alignment. It never runs $\widetilde g_{c,b}$ on a task and does not prove cost reduction, correctness, generalization, or preservation of task utility. Those properties require the later deployment measurements specified by the evaluation objective in Section~\ref{sec:problem}.

\subsection{Cycle completion and handoff}

After every eligible record has been processed or explicitly excluded for oracle contamination, the final fold is $g_c^\star=g_{c,B_c}$. A completed identity leaves the deployed harness unchanged. A completed change refreshes the stable tool runtime and makes $g_c^\star$ eligible to drive the next rollout, where its effects can be measured and can generate evidence for another cycle. The implementation records the rollout provenance, annotation state, proposal records, batch prompts, diffs, decisions, and a snapshot of the completed harness.

The current orchestrator can resume several completed stages, but a caught stage failure stops the cycle and in-place interruption is not a transaction-safe handoff. A failure-complete experiment must therefore place the cycle inside the external wrapper defined in Section~\ref{sec:problem}, which retains $h_c$ whenever the cycle does not finish with a completed change. This separation prevents an incomplete adaptation attempt from being counted as a successful or zero-cost update.
